\section{Coordenadas Polares}

Um sistema de \dt{coordenadas polares} é estabelecido da seguinte forma:

\begin{enumerate}
	\item Fixamos um ponto no plano, conhecido como \dt{polo} (ou origem), denotado por $O$.
	
	\item Fixamos uma semi-reta a partir de $O$, chamada \dt{eixo polar}.
\end{enumerate}

As \dt{coordenadas polares} de ponto {\color{red}$P$} do plano são formadas por um par ordenado $(r,\theta)$, sendo $r$ a distância de {\color{red}$P$} a $O$ e $\theta$ o ângulo entre o eixo polar e a reta $OP$.
\medskip


\begin{center}
	\begin{tikzpicture}[scale=3]
	\tikzset{>=latex}
	
	\coordinate (A) at ({2*cos(45)},{2*sin(45)});
	\coordinate (O) at (0,0);
	\coordinate (i) at (2,0);
	
	\pic[draw,angle eccentricity=1.5,->] {angle = i--O--A};
	
	\draw[->] (0,0) -- (3,0) node[right] {eixo polar};
	
	\node[thick] at (0.8,0.3) {$\theta$};
	\draw[very thick,blue] (0,0) -- ({2*cos(45)},{2*sin(45)});
	
	\node at (0.5,1) {$r$};
	\node[right] at (A) {$(r,\theta)$};
	\fill[red] (A) circle (1pt);
	
	\node[left] at (0,0) {$O$};
	
\end{tikzpicture}
\end{center}

{\large\textbf{Convenções}}
\begin{itemize}
	\item Um ângulo é \dt{positivo} se for medido no sentido anti-horário a partir do eixo polar e {\color{red}negativo} se for medido no sentido horário.
	
	\item $(0,\theta)$ representa o {\color{blue}polo} para qualquer valor de $\theta$.
	
	\item Se $r<0$, a coordenada $(r,\theta)$ corresponde ao simétrico, em relação à origem, do ponto $(|r|,\theta)$.
	
	
	\centering
	
	\begin{tikzpicture}
		\tikzset{>=latex}
		
		\coordinate (A) at ({2*cos(45)},{2*sin(45)});
		\coordinate (O) at (0,0);
		\coordinate (i) at (2,0);
		\coordinate (B) at ({2*cos(45+180)},{2*sin(45+180)});
		
		\pic[draw,->,angle radius=1cm] {angle = i--O--A};
		
		\pic[draw,->] {angle = i--O--B};
		
		\draw[->] (0,0) -- (3,0) node[right] {eixo polar};
		
		\node[thick] at (1.2,0.3) {$\theta$};
		\draw[very thick,dashed] (0,0) -- ({2*cos(45)},{2*sin(45)});
		
		\node at (0.5,1) {$|r|$};
		\node[right] at (A) {$(|r|,\theta)$};
		\fill[black] (A) circle (2pt);
		
		\node[below] at (0,0) {$O$};
		
		\draw[very thick,blue] (0,0) -- ({2*cos(45+180)},{2*sin(45+180)});
		
		\node[left] at (B) {$(r,\theta)$};
		\fill[red] (B) circle (2pt);
		\node at ({0.8*cos(90+45)},{0.8*sin(90+45)}) {$\theta +\pi$};
		
	\end{tikzpicture}
\end{itemize}